\documentclass[12pt]{article}
\usepackage[utf8]{inputenc}
\usepackage[T1]{fontenc}
\usepackage{hyperref}
\title{Tairuma Language Resources for K-12 Curriculum Development}
\author{Compiled by Claude Sonnet 5 (claude-sonnet-5) for Andr\'as Kornai}
\date{}
\begin{document}
\maketitle
\section{Introduction}
Tairuma (ISO 639-3 uar; also called Uaripi) is an Eleman language of Gulf Province, Papua New Guinea, in Central Kerema Rural LLG. Wikipedia cites about 4,500 speakers (2004), and Joshua Project gives 8,400 for the people group, which it calls Tiaru Ma'a. Sources report that Scripture exists only in portions and that no Tairuma textbooks exist, although early grades are said to be taught in the mother tongue.
\section{Primary Educational Resources}
No primers or textbooks were found in the searches.
\section{Linguistic and Reference Materials}
\begin{itemize}
\item Wikipedia, Tairuma language: \url{https://en.wikipedia.org/wiki/Tairuma_language}. It mentions a tentative grammar manuscript (Ikamu and Jo, 2014) and an SIL sociolinguistic profile (Rueck, Potter and Vila, 2010). I did not find direct links to either.
\item Glottolog, Tairuma: \url{https://glottolog.org/resource/languoid/id/tair1256}
\item Brown 1973, Eleman family, Trans New Guinea database entry: \url{https://transnewguinea.org/source/brown-1973}
\item Numeral system entry (MPI-EVA): \url{https://lingweb.eva.mpg.de/channumerals/Tairuma.htm}
\end{itemize}
\section{Religious and Community Resources}
\begin{itemize}
\item Joshua Project, Tiaru Ma'a in Papua New Guinea: \url{https://joshuaproject.net/people_groups/15676}
\end{itemize}
\section{Digital and Audio Resources}
None found.
\section{Curriculum Development Context}
The language is vigorous, but written resources are scarce. Curriculum work would need new primers and a practical orthography, building on the grammar manuscript and sociolinguistic profile once obtained from their authors or SIL.
\end{document}